\begin{pastexam}[label={px:ccu-108-csb-01}]{中正資工所乙組 108 學年度 第 1 題}
（問答）

(20pts) Quicksort is a basic sorting algorithm. (a) Please describe how it works. (b) Please
analyze its time complexity in the worst case and in the average case.
\end{pastexam}

\begin{pastsoln}
\textbf{(a) 做法}。 快速排序是分治法（資結書第~14.5~節）。 要排 $A[p..r]$：
\begin{enumerate}
\item \textbf{選基準}：從這一段挑一個數 $x$ 當基準（pivot），例如最右邊的 $A[r]$。
\item \textbf{分割}：把這一段重新排列成「$\le x$ 的數、$x$、$> x$ 的數」三部分，$x$ 落在位置
$q$。 分割完 $x$ 已經在最終的位置，以後不用再動。
\item \textbf{遞迴}：對左邊 $A[p..q-1]$ 和右邊 $A[q+1..r]$ 分別遞迴排序。 兩邊各自排好，整段就
排好了，不需要合併的步驟。
\end{enumerate}
只剩 $0$ 或 $1$ 個數的段（$p \ge r$）已經排好，直接返回。 用資結書虛擬碼~14.8、14.9 的 Lomuto
分割寫出來：

\begin{pseudo}{快速排序（Lomuto 分割）}
\Procedure{Quicksort}{$A, p, r$}
    \If{$p < r$}
        \State $q \gets$ \Call{Partition}{$A, p, r$}
        \State \Call{Quicksort}{$A, p, q - 1$}
        \State \Call{Quicksort}{$A, q + 1, r$}
    \EndIf
\EndProcedure
\Function{Partition}{$A, p, r$}
    \State $x \gets A[r]$ \Comment{基準}
    \State $i \gets p$ \Comment{$A[p..i-1]$ 都 $\le x$}
    \For{$j \gets p$ \To $r - 1$}
        \If{$A[j] \le x$}
            \State 交換 $A[i]$ 和 $A[j]$
            \State $i \gets i + 1$
        \EndIf
    \EndFor
    \State 交換 $A[i]$ 和 $A[r]$ \Comment{基準放到中間}
    \State \Return $i$
\EndFunction
\end{pseudo}

分割的迴圈維持「$A[p..i-1]$ 都 $\le x$，$A[i..j-1]$ 都 $> x$」：$A[j] \le x$ 時把它換到 $\le x$
那一區的尾端，否則留在原地，$> x$ 那一區自然變長。 迴圈結束時把基準和 $A[i]$（$> x$ 區的第一個）
交換，基準就在兩區中間。 例如 $\langle 3, 8, 2, 5\rangle$ 以 $5$ 為基準，分割成
$\langle 3, 2, \mathbf{5}, 8\rangle$，再分別排 $\langle 3, 2\rangle$ 和 $\langle 8\rangle$。

\textbf{(b) 時間複雜度}。 分割一段長度 $m$ 的陣列，迴圈做 $m - 1$ 次比較，花 $\Theta(m)$。 令
$T(n)$ 是排 $n$ 個數的時間，基準最後在第 $k$ 小的位置（$1 \le k \le n$）時
\[
T(n) = T(k - 1) + T(n - k) + \Theta(n).
\]

\textbf{最壞情況 $\Theta(n^2)$}。 每次分割都最不平均，基準是最大或最小值，切成 $n - 1$ 和 $0$：
\[
T(n) = T(n - 1) + T(0) + \Theta(n) = T(n-1) + \Theta(n)
= \Theta\Bigl(\sum_{m=1}^{n} m\Bigr) = \Theta(n^2).
\]
比較次數是 $(n-1) + (n-2) + \cdots + 1 = n(n-1)/2$。 取最右邊當基準時，已經排好（或反向排好）的
輸入就是這種情況。 這也是最壞：每一層分割至少讓一個數（基準）不再參加，遞迴深度最多 $n$，
每一層的比較總數不超過 $n$，所以不會比 $O(n^2)$ 更差。

\textbf{最好情況 $\Theta(n \lg n)$}。 每次都切成兩半時 $T(n) = 2T(n/2) + \Theta(n)$，由主定理
（資結書第~2.4~節）得 $\Theta(n \lg n)$。

\textbf{平均情況 $\Theta(n \lg n)$}。 假設 $n$ 個數互不相同，$n!$ 種排列出現的機會相同（或基準是
隨機選的）。 基準是第 $k$ 小的機率都是 $1/n$，而且分割完左右兩段仍然是隨機排列。 令 $C(n)$ 是
平均比較次數，$C(0) = C(1) = 0$：
\[
C(n) = (n - 1) + \frac{1}{n} \sum_{k=1}^{n} \bigl(C(k-1) + C(n-k)\bigr)
     = (n - 1) + \frac{2}{n} \sum_{k=0}^{n-1} C(k).
\]
兩邊乘 $n$，再減去 $n - 1$ 時的同一個式子，消掉總和：
\[
nC(n) - (n-1)C(n-1) = 2(n-1) + 2C(n-1),
\qquad\text{即}\qquad
nC(n) = (n+1)C(n-1) + 2(n-1).
\]
兩邊除以 $n(n+1)$，並用 $\dfrac{2(n-1)}{n(n+1)} = \dfrac{4}{n+1} - \dfrac{2}{n}$：
\[
\frac{C(n)}{n+1} = \frac{C(n-1)}{n} + \frac{4}{n+1} - \frac{2}{n}.
\]
從 $\frac{C(1)}{2} = 0$ 一路加到 $n$，令 $H_n = 1 + \frac12 + \cdots + \frac1n$：
\[
\frac{C(n)}{n+1} = \sum_{m=2}^{n}\Bigl(\frac{4}{m+1} - \frac{2}{m}\Bigr)
= 4\Bigl(H_{n+1} - \frac32\Bigr) - 2(H_n - 1)
= 2H_n + \frac{4}{n+1} - 4,
\]
所以
\[
C(n) = 2(n+1)H_n - 4n \approx 2n \ln n \approx 1.39\, n \lg n,
\]
平均時間 $\Theta(n \lg n)$，和資結書定理~14.3 的結果相同（那裡用「任兩個數比較的機率」另外推了
一次）。 例如 $n = 4$：$C(4) = 2 \cdot 5 \cdot \frac{25}{12} - 16 = \frac{29}{6}$，$24$ 種排列
的比較次數平均起來正是 $29/6$。

總結：最壞 $\Theta(n^2)$，平均 $\Theta(n \lg n)$（最好也是 $\Theta(n \lg n)$）。 實務上用隨機選
基準或三數取中，讓最壞情況不會被特定的輸入（例如已排好的資料）觸發。
\end{pastsoln}
